% Phase VII section -- insert after Phase VI.
% Requires: \usepackage{amsmath}, \usepackage{booktabs}, \usepackage{tabularx},
%           \usepackage{float}, \usepackage{hyperref}

\section{Phase VII: research application, reproducibility, and deployment}
\label{sec:phase7}

\subsection{Objective and scientific boundary}

Phase VII converts the locked Phase VI experiment into an executable research
demonstrator. It does not train, tune, or select another predictive model. Its purpose is
to verify that the selected computer-vision pipeline can be exposed through a documented
API, consumed by a web interface, tested automatically, packaged reproducibly, and prepared
for cloud deployment. This separation prevents application engineering from reopening the
scientific model-selection protocol.

The deployed contract is the Phase VI configuration selected on Scene18 before the one-time
Scene20 evaluation: a frozen DINOv2 ViT-S/14 representation, a multinomial logistic probe on
the semantic-sky region, and causal exponential moving-average smoothing with
$\alpha=0.6$. The output classes are \texttt{clone}, \texttt{fog}, and \texttt{rain}.
Because the regional input relies on the Virtual KITTI 2 ground-truth colour mask, the
application is explicitly presented as a synthetic-data research demonstrator rather than
as a validated real-road weather or safety system.

\subsection{System architecture}

The system separates inference, presentation, automation, and deployment. A FastAPI
backend validates the uploaded RGB frame and its corresponding semantic mask, reconstructs
the semantic-sky image used in Phase VI, extracts the frozen DINOv2 embedding, applies the
serialized logistic-regression model, and updates a causal EMA state identified by the
sequence identifier. A React and TypeScript interface uploads ordered frames and displays
raw and temporally smoothed probabilities, confidence, entropy, abstention status, and
latency.

\begin{table}[H]
\centering
\small
\caption{Phase VII application components.}
\label{tab:phase7_components}
\begin{tabularx}{\textwidth}{p{3.2cm}p{3.0cm}X}
\toprule
Component & Technology & Role\\
\midrule
Inference service & FastAPI, Python & Typed HTTP interface, input validation, DINOv2 inference, probability output, and sequence-state management\\
Web interface & React, TypeScript, Vite & Upload of RGB--mask pairs and interactive display of temporal predictions\\
Packaging & Docker, Docker Compose & Reproducible API and interface environments for local execution\\
Verification & pytest, Ruff, TypeScript compiler & Unit, endpoint, style, and production-build checks\\
Automation & GitHub Actions & Continuous integration and immutable image deployment using the Git commit SHA\\
Cloud target & Azure Container Apps, ACR & Managed container revisions, ingress, health probes, logs, and operational metrics\\
Artifact source & Google Cloud Storage & Retrieval of the checksummed canonical Phase VI release before the API image is built\\
\bottomrule
\end{tabularx}
\end{table}

\subsection{Inference API and temporal state}

For frame $t$ of sequence $s$, the locked classifier returns a probability vector
$\mathbf{p}_{s,t}$. The API applies exactly the Phase VI recursion

\begin{equation}
\widetilde{\mathbf{p}}_{s,t}=
\begin{cases}
\mathbf{p}_{s,t}, & t=0,\\
0.6\,\mathbf{p}_{s,t}+0.4\,\widetilde{\mathbf{p}}_{s,t-1}, & t>0.
\end{cases}
\end{equation}

State is isolated by sequence identifier and can be reset explicitly. It is bounded by a
maximum number of sequences and an inactivity timeout. The response reports the predicted
class, smoothed confidence, entropy, raw and smoothed probability vectors, frame index, and
measured inference time. An optional confidence threshold supports selective prediction:
frames below the threshold are flagged as abstentions rather than silently treated as
reliable decisions.

The service exposes separate liveness and readiness endpoints. Liveness confirms that the
API process responds, whereas readiness confirms that the Phase VI artifacts and DINOv2
model have loaded successfully. If the canonical release manifest is available, the model
and region-definition files are checked against their SHA-256 hashes before serving
predictions.

\subsection{Automated verification and continuous integration}

Backend tests verify that EMA is causal, sequence-local, resettable, and equal to the
specified recursion. Endpoint tests verify health reporting, model metadata, valid
multipart inference, temporal updates, and rejection of invalid image payloads. The
continuous-integration workflow runs these tests and Ruff static checks with Python 3.12,
then installs and compiles the React/TypeScript production bundle with Node.js 24.

A deterministic mock engine is included only for application smoke tests when the private
Phase VI artifacts are unavailable. It is separated through the \texttt{MODEL\_MODE}
configuration and must not be used to report scientific performance. All thesis metrics
remain those of the locked Phase VI release.

\subsection{Containerisation and cloud-deployment design}

Docker Compose connects the API and interface locally. For the cloud workflow, GitHub
Actions authenticates to Google Cloud Storage, downloads the locked Phase VI archive,
verifies the required artifact names, builds images tagged with the immutable commit SHA,
and deploys two revisions to Azure Container Apps through Azure Container Registry. The
frontend is built with the deployed API URL, and the API permits only the resulting frontend
origin.

The application provides Prometheus-format request and latency metrics and emits structured
container logs. Recommended monitoring covers error ratio, readiness failures, restart
loops, p95 inference latency, abstention rate, class-frequency drift, and confidence drift.
Uploaded images are processed in memory and are not persisted by the application.

The current EMA store is process-local. Consequently, the demonstrator is deliberately
limited to one API replica so that consecutive frames of a sequence cannot be routed to
different states. Production horizontal scaling would require a shared state store or a
stateless protocol that carries temporal state between requests.

\subsection{Outcome and limitations}

Phase VII demonstrates the engineering path from a locked notebook experiment to a typed,
testable, containerised application while maintaining the model identity and scientific
limitations established in Phase VI. It adds no evidence that the very high Virtual KITTI 2
Scene20 score transfers to real traffic. The strongest remaining limitation is the oracle
semantic-sky mask. Replacing it with predicted segmentation, measuring end-to-end latency,
and evaluating independent real adverse-weather videos are prerequisites for any claim of
operational generalisation. An actual Azure production deployment additionally depends on
user-controlled cloud identities and is not claimed solely from the presence of deployment
configuration files.
